\subsection{Zielbild des Technologie-Templates}
\label{subsec:zielbild-technologie-template}

Das Template soll eine wiederverwendbare technische Ausgangsbasis für Web-,
Cloud- und Desktop-Projekte bilden. Wiederverwendung setzt dabei geplante
Gemeinsamkeiten und beherrschte Variabilität voraus
\parencite{krueger1992software-reuse,clements2001software-product-lines}.
Projektstruktur, Entwicklungsworkflow, Tooling, Dokumentationsprinzipien,
Konfigurationslogik, Tests und Schnittstellen sollen deshalb in einem
gemeinsamen Kern zentral gepflegt werden.

Nicht jedes Projekt benötigt sämtliche Komponenten. Profile sollen passende
Kombinationen des Kerns auswählen; optionale Capabilities wie PostgreSQL
sollen kompatible Backend-Profile unabhängig vom Plattformprofil erweitern
\parencite{kleiveist2026profiles}. Bereits im Zielbild sind Frontend, Backend,
Desktop-Schicht, Persistenz, Tooling und Deployment als getrennte
Verantwortungsbereiche vorgesehen. Ob die Umsetzung diese Anforderungen
erfüllt, wird erst in den nachfolgenden Architektur-, Verifikations- und
Bewertungskapiteln geprüft.
